\chapter{The Ricci Calculus}

\section*{A.1. Component representations of tensor fields}

\begin{definition}[Component representation of a tensor field]
\label{appa-def-component-representation}
\source{chow-knopf-ricci-flow:page-0292}
\dcref{Appendix A.1}
\group{component-representations-of-tensor-fields}

Let \((\mathcal M^n,g)\) be a smooth Riemannian manifold, let \(U\subset \mathcal M^n\) be open, and let \(\varphi:U\to \mathbb R^n\) be a chart with induced local coordinates \((x^1,\ldots,x^n)\). For a \((p,q)\)-tensor field \(A\), its component representation with respect to \(\varphi\) is the tensor field \(\bigl(A^{k_1\cdots k_q}_{j_1\cdots j_p}\bigr)\) defined by
\[
A = A^{k_1\cdots k_q}_{j_1\cdots j_p}\, dx^{j_1}\otimes \cdots \otimes dx^{j_p}\otimes \frac{\partial}{\partial x^{k_1}}\otimes \cdots \otimes \frac{\partial}{\partial x^{k_q}}.
\]
Equivalently,
\[
A^{k_1\cdots k_q}_{j_1\cdots j_p}
=
A\!\left(
\frac{\partial}{\partial x^{j_1}},\ldots,\frac{\partial}{\partial x^{j_p}};\,
dx^{k_1},\ldots,dx^{k_q}
\right).
\]
\end{definition}

\section*{A.2. First-order differential operators on tensors}

\begin{definition}[Metric dual of a vector field]
\label{appa-def-metric-dual-vector}
\source{chow-knopf-ricci-flow:page-0293}
\dcref{Section A.2}
\group{first-order-differential-operators-on-tensors}

For a vector field \(X\in C^\infty(T\mathcal M^n)\), its metric dual is the covector field \(X_b\in C^\infty(T^*\mathcal M^n)\) defined in local coordinates by
\[
X_b=(X^i g_{ij})\,dx^j.
\]
\end{definition}

\begin{definition}[Metric dual of a covector field]
\label{appa-def-metric-dual-covector}
\source{chow-knopf-ricci-flow:page-0293}
\dcref{Section A.2}
\group{first-order-differential-operators-on-tensors}

For a covector field \(\theta\in C^\infty(T^*\mathcal M^n)\), its metric dual is the vector field \(\theta^\sharp\in C^\infty(T\mathcal M^n)\) defined in local coordinates by
\[
\theta^\sharp = (\theta_i g^{ij})\,\frac{\partial}{\partial x^j}.
\]
\end{definition}

\begin{remark}[Index raising and lowering]
\label{appa-rem-index-raising-lowering}
\source{chow-knopf-ricci-flow:page-0293}
\dcref{Section A.2}
\group{first-order-differential-operators-on-tensors}

The metric and its inverse induce canonical isomorphisms between tensor bundles of the same total valence. In local coordinate formulas, the source freely writes quantities such as \(R^k{}_i=g^{jk}R_{ij}\) and \(R^{k\ell}=g^{ik}g^{j\ell}R_{ij}\).
\end{remark}

\begin{definition}[Action of a vector field on a function]
\label{appa-def-vector-action-on-function}
\source{chow-knopf-ricci-flow:page-0293}
\dcref{Section A.2}
\group{first-order-differential-operators-on-tensors}

If \(X\in C^\infty(T\mathcal M^n)\) and \(f:\mathcal M^n\to \mathbb R\) is differentiable, then \(X\) acts on \(f\) by differentiation, with local expression
\[
X(f)=X^i\frac{\partial f}{\partial x^i}.
\]
\end{definition}

\begin{definition}[Levi-Civita covariant derivative]
\label{appa-def-levi-civita-covariant-derivative}
\source{chow-knopf-ricci-flow:page-0293}
\uses{appa-def-vector-action-on-function}
\dcref{Section A.2}
\group{first-order-differential-operators-on-tensors}


The Levi-Civita connection of \(g\) is identified with the covariant derivative, the unique operator
\[
\nabla : C^\infty(T\mathcal M^n)\to C^\infty(T\mathcal M^n\otimes T^*\mathcal M^n)
\]
such that
\[
X\bigl(g(Y,Z)\bigr)=g(\nabla_XY,Z)+g(Y,\nabla_XZ)
\]
for all vector fields \(X,Y,Z\).
\end{definition}

\begin{remark}[Covariant derivative with respect to a vector]
\label{appa-rem-covariant-derivative-vector}
\source{chow-knopf-ricci-flow:page-0293}
\dcref{Section A.2}
\group{first-order-differential-operators-on-tensors}

For vector fields \(X,Y\), the notation \(\nabla_YX=(\nabla X)(Y)\) is used.
\end{remark}

\begin{definition}[Covariant derivative of a vector field]
\label{appa-def-covariant-derivative-vector}
\lean{ChowKnopf.covariantDerivativeVectorComponent,
  ChowKnopf.covariantDerivativeVectorComponent_eq}
\leanok
\source{chow-knopf-ricci-flow:page-0294}
\dcref{Section A.2}
\group{first-order-differential-operators-on-tensors}

In a local coordinate chart, the covariant derivative of a vector field \(X\) has components
\[
\nabla_i X^k = \frac{\partial}{\partial x^i}X^k + \Gamma^k_{ij}X^j,
\]
and one customarily writes \(\nabla_{\partial/\partial x^i}\) as \(\nabla_i\).
\end{definition}

\begin{definition}[Covariant derivative of a covector field]
\label{appa-def-covariant-derivative-covector}
\lean{ChowKnopf.LeviCivitaConnection, ChowKnopf.covariantDerivativeCovector,
  ChowKnopf.covariantDerivativeCovector_characterization}
\leanok
\source{chow-knopf-ricci-flow:page-0294}
\dcref{Section A.2}
\group{first-order-differential-operators-on-tensors}

The Levi-Civita connection induces a covariant derivative \(\nabla^*\) on covector fields, characterized by
\[
X\bigl(\theta(Y)\bigr) = (\nabla_X^*\theta)(Y) + \theta(\nabla_X Y)
\]
for all vector fields \(X,Y\) and covector fields \(\theta\).
\end{definition}

\begin{definition}[Covariant derivative on tensor bundles]
\label{appa-def-covariant-derivative-general-tensor}
\lean{ChowKnopf.covariantDerivativeTensorEvaluation,
  ChowKnopf.covariantDerivativeTensorEvaluation_characterization}
\source{chow-knopf-ricci-flow:page-0294}
\uses{appa-def-covariant-derivative-vector, appa-def-covariant-derivative-covector}
\dcref{Section A.2}
\group{first-order-differential-operators-on-tensors}


For each \((p,q)\)-tensor bundle, the Levi-Civita connection determines a covariant derivative
\[
\nabla^{(p,q)} : C^\infty\!\left(T_p^q \mathcal M^n\right) \to C^\infty\!\left(T_{p+1}^q \mathcal M^n\right)
\]
defined by the rule
\[
X\bigl(A(Y_1,\ldots,Y_p;\,\theta_1,\ldots,\theta_q)\bigr)
= \bigl(\nabla_X^{(p,q)}A\bigr)(Y_1,\ldots,Y_p;\,\theta_1,\ldots,\theta_q)
\]
\[
\qquad + \sum_{i=1}^p A(Y_1,\ldots,\nabla_X Y_i,\ldots,Y_p;\,\theta_1,\ldots,\theta_q)
\]
\[
\qquad + \sum_{j=1}^q A(Y_1,\ldots,Y_p;\,\theta_1,\ldots,\nabla_X^*\theta_j,\ldots,\theta_q).
\]
The source notes that \(\nabla^{(0,1)}=\nabla\) and \(\nabla^{(1,0)}=\nabla^*\).
\end{definition}

\begin{definition}[Component representation of the covariant derivative]
\label{appa-def-component-representation-covariant-derivative}
\source{chow-knopf-ricci-flow:page-0294}
\uses{appa-def-component-representation, appa-def-covariant-derivative-general-tensor}
\dcref{Section A.2}
\group{first-order-differential-operators-on-tensors}


For a \((p,q)\)-tensor field \(A\), the component representation of its covariant derivative is defined by
\[
\nabla_i^{(p,q)} A^{k_1\cdots k_q}_{j_1\cdots j_p}
\doteq
(\nabla A)\!\left(\frac{\partial}{\partial x^i}, \frac{\partial}{\partial x^{j_1}}, \ldots, \frac{\partial}{\partial x^{j_p}};\, dx^{k_1}, \ldots, dx^{k_q}\right).
\]
\end{definition}

\begin{remark}[Abuse of notation for covariant derivatives]
\label{appa-rem-abuse-notation-covariant-derivatives}
\source{chow-knopf-ricci-flow:page-0294}
\dcref{Remark A.1}
\group{first-order-differential-operators-on-tensors}

It is customary to denote all operators \(\nabla^{(p,q)}\) simply by \(\nabla\).
\end{remark}

\begin{definition}[Lie derivative on vector fields]
\label{appa-def-lie-derivative-vector}
\source{chow-knopf-ricci-flow:page-0295}
\dcref{Section A.2}
\group{first-order-differential-operators-on-tensors}

The Lie derivative of vector fields is defined by
\[
\Lie_X Y=[X,Y]
\]
for all vector fields \(X,Y\), where \([X,Y]\) is the unique vector field satisfying
\[
[X,Y](f)=X(Y(f))-Y(X(f))
\]
for all differentiable functions \(f:\mathcal M^n\to\mathbb R\).
\end{definition}

\begin{definition}[Lie derivative on covector fields]
\label{appa-def-lie-derivative-covector}
\source{chow-knopf-ricci-flow:page-0295}
\dcref{Section A.2}
\group{first-order-differential-operators-on-tensors}

If \(\theta\) is a covector field and \(X,Y\) are vector fields, then the Lie derivative of \(\theta\) is given by
\[
(\Lie_X \theta)(Y)=X(\theta(Y))-\theta([X,Y]).
\]
Equivalently,
\[
(\Lie_X \theta)(Y)=d\theta(X,Y)+Y(\theta(X)),
\]
where \(d\) is the exterior derivative.
\end{definition}

\begin{remark}[Lie derivative of a vector field via the Levi-Civita connection]
\label{appa-rem-lie-bracket-via-connection}
\source{chow-knopf-ricci-flow:page-0295}
\dcref{Equation A.1}
\group{first-order-differential-operators-on-tensors}

The Lie derivative of a vector field may be computed from the Levi-Civita connection by
\[
\Lie_XY=[X,Y]=\nabla_XY-\nabla_YX.
\]
\end{remark}

\begin{remark}[Lie derivative of covectors and tensors]
\label{appa-rem-lie-derivative-tensors}
\source{chow-knopf-ricci-flow:page-0295}
\dcref{Equation A.2}
\group{first-order-differential-operators-on-tensors}

For a covector field \(\theta\),
\[
(\Lie_X \theta)(Y)=(\nabla_X\theta)(Y)+\theta(\nabla_YX).
\]
More generally, for a \((p,q)\)-tensor field \(A\),
\[
(\Lie_X A)(Y_1,\ldots,Y_p;\theta_1,\ldots,\theta_q)
=X\bigl(A(Y_1,\ldots,Y_p;\theta_1,\ldots,\theta_q)\bigr)
\]
\[
\qquad - \sum_{1\le i\le p} A(Y_1,\ldots,[X,Y_i],\ldots,Y_p;\theta_1,\ldots,\theta_q)
\]
\[
\qquad - \sum_{1\le j\le q} A(Y_1,\ldots,Y_p;\theta_1,\ldots,\Lie_X\theta_j,\ldots,\theta_q).
\]
\end{remark}

\begin{remark}[Lie derivative of the metric]
\label{appa-rem-lie-derivative-metric}
\source{chow-knopf-ricci-flow:page-0295}
\dcref{Section A.2}
\group{first-order-differential-operators-on-tensors}

The Lie derivative of the metric satisfies
\[
(\Lie_X g)(Y,Z)=g(\nabla_YX,Z)+g(Y,\nabla_ZX),
\]
and in local coordinates
\[
(\Lie_X g)_{ij}=\nabla_i X_j+\nabla_j X_i.
\]
\end{remark}

\begin{definition}[Divergence operator on symmetric tensors]
\label{appa-def-divergence-operator}
\source{chow-knopf-ricci-flow:page-0295, chow-knopf-ricci-flow:page-0296}
\uses{appa-def-lie-derivative-vector}
\dcref{Section A.2}
\group{first-order-differential-operators-on-tensors}


The divergence is the operator
\[
\delta : C^\infty(T^*\mathcal M^n \otimes_S T^*\mathcal M^n)\to C^\infty(T^*\mathcal M^n)
\]
defined for every symmetric tensor \(A\) and every vector field \(X\) by
\[
(\delta A)(X)=-(\divergence A)(X)=\sum_{i=1}^n (\nabla A)(e_i,e_i,X),
\]
where \(\{e_i\}_{i=1}^n\) is a local orthonormal frame field.
\end{definition}

\begin{definition}[Formal adjoint of the divergence]
\label{appa-def-formal-adjoint-divergence}
\source{chow-knopf-ricci-flow:page-0296}
\dcref{Section A.2}
\group{first-order-differential-operators-on-tensors}

If \(\mathcal M^n\) is compact, the formal adjoint of \(\delta\) with respect to the \(L^2\) inner product
\[
(\cdot,\cdot)=\int_{\mathcal M^n}\langle \cdot,\cdot\rangle\,d\mu
\]
is denoted by
\[
\delta^* : C^\infty(T^*\mathcal M^n)\to C^\infty(T^*\mathcal M^n\otimes_S T^*\mathcal M^n).
\]
It satisfies
\[
\delta^*\theta=\frac12 \Lie_{\theta^\sharp}g
\]
for all covector fields \(\theta\).
\end{definition}

\begin{remark}[Extending \(\delta^*\) beyond the compact case]
\label{appa-rem-extending-divergence-adjoint}
\source{chow-knopf-ricci-flow:page-0296}
\dcref{Section A.2}
\group{first-order-differential-operators-on-tensors}

The formula \(\delta^*\theta=\tfrac12 \Lie_{\theta^\sharp}g\) can be used to define \(\delta^*\) even when \(\mathcal M^n\) is not compact.
\end{remark}

\section*{A.3. First-order differential operators on forms}

\begin{definition}[Exterior derivative on forms]
\label{appa-def-exterior-derivative}
\source{chow-knopf-ricci-flow:page-0296, chow-knopf-ricci-flow:page-0297}
\dcref{Section A.3}
\group{first-order-differential-operators-on-forms}

Let \(\Omega^p(\mathcal M^n)\coloneqq C^\infty\bigl(\Lambda^p(T^*\mathcal M^n)\bigr)\). The exterior derivative is the family of operators
\[
d\equiv d_p:\Omega^p(\mathcal M^n)\to \Omega^{p+1}(\mathcal M^n)
\]
defined for all \(p\)-forms \(\theta\) and vector fields \(Y_0,\ldots,Y_p\) by
\[
d\theta(Y_0,\ldots,Y_p)
\coloneqq \sum_{0\le i\le p}(-1)^i\theta(Y_0,\ldots,\widehat{Y_i},\ldots,Y_p)
\]
\[
\qquad + \sum_{0\le i<j\le p}(-1)^{i+j}\theta([Y_i,Y_j],Y_0,\ldots,\widehat{Y_i},\ldots,\widehat{Y_j},\ldots,Y_p).
\]
\end{definition}

\begin{remark}[Exterior derivative in covariant form]
\label{appa-rem-exterior-derivative-covariant-form}
\source{chow-knopf-ricci-flow:page-0297}
\dcref{Section A.3}
\group{first-order-differential-operators-on-forms}

Although independent of the metric, the exterior derivative may be written in terms of the Levi-Civita covariant derivative as
\[
d\theta(Y_0,\ldots,Y_p)=\sum_{i=0}^p (-1)^i (\nabla_{Y_i}\theta)(Y_0,\ldots,\widehat{Y_i},\ldots,Y_p).
\]
\end{remark}

\begin{definition}[Formal adjoint of the exterior derivative]
\label{appa-def-formal-adjoint-exterior-derivative}
\source{chow-knopf-ricci-flow:page-0297}
\dcref{Section A.3}
\group{first-order-differential-operators-on-forms}

If \(\mathcal M^n\) is compact, the formal adjoint of \(d\) with respect to the \(L^2\) inner product is the family of operators
\[
\delta\equiv \delta_{p+1}:\Omega^{p+1}(\mathcal M^n)\to \Omega^p(\mathcal M^n)
\]
defined by
\[
(d\theta,\eta)=(\theta,\delta\eta)
\]
for all \(\theta\in \Omega^p(\mathcal M^n)\) and \(\eta\in \Omega^{p+1}(\mathcal M^n)\).
\end{definition}

\begin{remark}[Covariant formula for the adjoint of \(d\)]
\label{appa-rem-adjoint-exterior-derivative-covariant-form}
\source{chow-knopf-ricci-flow:page-0297}
\dcref{Section A.3}
\group{first-order-differential-operators-on-forms}

By integration by parts,
\[
(\delta\eta)(X_1,\ldots,X_p)
= - \sum_{i=1}^n (\nabla\eta)(e_i,e_i,X_1,\ldots,X_p),
\]
where \(\{e_i\}_{i=1}^n\) is a local orthonormal frame field. This formula also serves to define \(\delta\) when \(\mathcal M^n\) is not compact.
\end{remark}

\section*{A.4. Second-order differential operators}

\begin{definition}[Rough Laplacian]
\label{appa-def-rough-laplacian}
\source{chow-knopf-ricci-flow:page-0297}
\uses{appa-def-covariant-derivative-general-tensor}
\dcref{Section A.4}
\group{second-order-differential-operators}


The rough Laplacian is the family of operators
\[
\Delta : C^\infty(T_p^q \mathcal M^n)\to C^\infty(T_p^q \mathcal M^n)
\]
defined by
\[
(\Delta A)(Y_1,\ldots,Y_p;\theta_1,\ldots,\theta_q)
= \sum_{i=1}^n (\nabla^2 A)(e_i,e_i,Y_1,\ldots,Y_p;\theta_1,\ldots,\theta_q).
\]
\end{definition}

\begin{definition}[Component representation of the rough Laplacian]
\label{appa-def-rough-laplacian-components}
\source{chow-knopf-ricci-flow:page-0297}
\uses{appa-def-component-representation, appa-def-rough-laplacian}
\dcref{Section A.4}
\group{second-order-differential-operators}


In local coordinates, the component representation of the rough Laplacian is
\[
\Delta A^{k_1\cdots k_q}_{j_1\cdots j_p}
:= (\Delta A)\!\left(\frac{\partial}{\partial x^i},\frac{\partial}{\partial x^{j_1}},\ldots,\frac{\partial}{\partial x^{j_p}};dx^{k_1},\ldots,dx^{k_q}\right).
\]
\end{definition}

\begin{definition}[Hodge-de Rham Laplacian]
\label{appa-def-hodge-de-rham-laplacian}
\source{chow-knopf-ricci-flow:page-0297}
\uses{appa-def-exterior-derivative, appa-def-formal-adjoint-exterior-derivative}
\dcref{Section A.4}
\group{second-order-differential-operators}


The Hodge-de Rham Laplacian, also called the Laplace-Beltrami operator, is the family of maps
\[
-\Delta_d:\Omega^p(\mathcal M^n)\to \Omega^p(\mathcal M^n)
\]
defined by
\[
-\Delta_d=d\delta+\delta d=d_{p-1}\delta_p+\delta_{p+1}d_p.
\]
\end{definition}

\begin{remark}[One- and two-form components of the Hodge-de Rham Laplacian]
\label{appa-rem-hodge-laplacian-components}
\source{chow-knopf-ricci-flow:page-0297, chow-knopf-ricci-flow:page-0298}
\dcref{Section A.4}
\group{second-order-differential-operators}

For a \(1\)-form \(\theta\),
\[
\Delta_d\theta=(\Delta\theta_i-R_i^{\ j}\theta_j)\,dx^i.
\]
For a \(2\)-form \(\alpha\),
\[
\Delta_d\alpha
= \bigl(\Delta \alpha_{ij} + g^{kp}g^{\ell q}R_{ijk\ell}\alpha_{qp}
- g^{k\ell}R_{ik}\alpha_{\ell j}
- g^{k\ell}R_{jk}\alpha_{i\ell}\bigr)\,dx^i\wedge dx^j.
\]
\end{remark}

\begin{definition}[Lichnerowicz Laplacian on symmetric \(2\)-tensors]
\label{appa-def-lichnerowicz-laplacian}
\source{chow-knopf-ricci-flow:page-0298}
\uses{appa-def-rough-laplacian}
\dcref{Section A.4}
\group{second-order-differential-operators}


The Lichnerowicz Laplacian \(\Delta_L\) on symmetric \(2\)-tensors is the operator for which the source records the commutative diagram relating \(\Delta_L\), \(\delta\), \(\delta^*\), and \(\Delta_d\) when \((\mathcal M^n,g)\) is Ricci-parallel.
\end{definition}

\begin{remark}[Components of the Lichnerowicz Laplacian]
\label{appa-rem-lichnerowicz-laplacian-components}
\source{chow-knopf-ricci-flow:page-0298}
\dcref{Equation A.4}
\group{second-order-differential-operators}

If \(A\) is a symmetric \((2,0)\)-tensor, then
\[
\Delta_L A
= \left(\Delta A_{ij} + 2g^{kp}g^{\ell q}R_{ik\ell j}A_{pq} - g^{k\ell}R_{ik}A_{\ell j} - g^{k\ell}R_{jk}A_{\ell i}\right)\, dx^i \otimes dx^j .
\]
\end{remark}

\begin{remark}[Lichnerowicz and Hodge-de Rham comparison]
\label{appa-rem-lichnerowicz-hodge-comparison}
\source{chow-knopf-ricci-flow:page-0298}
\dcref{Remark A.2, Remark A.3}
\group{second-order-differential-operators}

The source notes that the Lichnerowicz Laplacian is essentially the linearization of the Ricci flow operator, and that on symmetric \(2\)-tensors it is formally the same as the Hodge-de Rham Laplacian on \(2\)-forms after applying the first Bianchi identity.
\end{remark}

\section*{A.5. Notation for higher derivatives}

\begin{definition}[Higher covariant derivatives]
\label{appa-def-higher-covariant-derivatives}
\source{chow-knopf-ricci-flow:page-0298}
\uses{appa-def-covariant-derivative-general-tensor}
\dcref{Section A.5}
\group{notation-for-higher-derivatives}


For \(k>1\), the source considers operators
\[
\nabla^k : C^\infty\!\left(T^q_p \mathcal M^n\right) \to C^\infty\!\left(T^q_{p+k} \mathcal M^n\right).
\]
If \(A\in C^\infty(T^q_p \mathcal M^n)\) and \(X_1,\ldots,X_k\) are vector fields, then
\[
\nabla_{X_1}\nabla_{X_2}\cdots\nabla_{X_k}A
\]
denotes the unique \((p,q)\)-tensor field satisfying
\[
(\nabla_{X_1}\nabla_{X_2}\cdots\nabla_{X_k}A)(Y_1,\ldots,Y_p;\,\theta_1,\ldots,\theta_q)
= (\nabla^k A)(X_1,X_2,\ldots,X_k,Y_1,\ldots,Y_p;\,\theta_1,\ldots,\theta_q).
\]
\end{definition}

\begin{remark}[Iterated covariant derivatives versus nested covariant differentiation]
\label{appa-rem-iterated-vs-nested-covariant-derivatives}
\source{chow-knopf-ricci-flow:page-0298}
\dcref{Section A.5}
\group{notation-for-higher-derivatives}

In the source's convention,
\[
(\nabla^k A)(X_1,\ldots,X_k):=\nabla_{X_1}\nabla_{X_2}\cdots\nabla_{X_k}A \neq \nabla_{X_1}(\nabla_{X_2}(\cdots(\nabla_{X_k}A)))
\]
in general. In particular,
\[
\nabla_X(\nabla_YA)=\nabla_X\nabla_YA+\nabla_{\nabla_XY}A=(\nabla^2A)(X,Y)+(\nabla A)(\nabla_XY).
\]
\end{remark}

\section*{A.6. Commuting covariant derivatives}

\begin{definition}[Riemann curvature tensor via second covariant derivatives]
\label{appa-def-riemann-curvature-via-covariant-derivatives}
\source{chow-knopf-ricci-flow:page-0299}
\uses{appa-def-higher-covariant-derivatives}
\dcref{Equation A.6}
\group{commuting-covariant-derivatives}


The \((3,1)\)-Riemann curvature tensor is defined by
\[
R(X,Y)Z=(\nabla^2 Z)(X,Y)-(\nabla^2 Z)(Y,X)
\]
for all vector fields \(X,Y,Z\).
\end{definition}

\begin{remark}[Curvature as a commutator of covariant derivatives]
\label{appa-rem-curvature-commutator}
\source{chow-knopf-ricci-flow:page-0299}
\dcref{Equation A.7}
\group{commuting-covariant-derivatives}

Following the source's convention for iterated covariant derivatives,
\[
R(X,Y)Z=\nabla_X\nabla_YZ-\nabla_Y\nabla_XZ.
\]
\end{remark}

\begin{remark}[Coordinate form of the curvature tensor]
\label{appa-rem-curvature-coordinate-form}
\source{chow-knopf-ricci-flow:page-0299}
\dcref{Section A.6}
\group{commuting-covariant-derivatives}

In local coordinates, the coordinate vector fields commute, so
\[
R\!\left(\frac{\partial}{\partial x^i},\frac{\partial}{\partial x^j}\right)\frac{\partial}{\partial x^k}
=R^\ell{}_{ijk}\frac{\partial}{\partial x^\ell}
=\left(\nabla_i\nabla_j-\nabla_j\nabla_i\right)\!\left(\frac{\partial}{\partial x^k}\right).
\]
\end{remark}

\begin{remark}[Ricci identities]
\label{appa-rem-ricci-identities}
\source{chow-knopf-ricci-flow:page-0299}
\dcref{Section A.6}
\group{commuting-covariant-derivatives}

The commutator formulas implied by the curvature tensor are called the Ricci identities. For a vector field \(X\),
\[
[\nabla_i,\nabla_j]X^\ell \equiv \nabla_i\nabla_jX^\ell-\nabla_j\nabla_iX^\ell = R^\ell{}_{ijk}X^k.
\]
For a covector field \(\theta\),
\[
[\nabla_i,\nabla_j]\theta_k \equiv \nabla_i\nabla_j\theta_k-\nabla_j\nabla_i\theta_k = -R^\ell{}_{ijk}\theta_\ell.
\]
\end{remark}

\begin{remark}[Commutator on general tensors]
\label{appa-rem-commutator-general-tensors}
\source{chow-knopf-ricci-flow:page-0299}
\dcref{Section A.6}
\group{commuting-covariant-derivatives}

If \(A\) is any \((p,q)\)-tensor field, then
\[
[\nabla_i,\nabla_j]A^{\ell_1\cdots \ell_q}{}_{k_1\cdots k_p}
\equiv \nabla_i\nabla_j A^{\ell_1\cdots \ell_q}{}_{k_1\cdots k_p}
-\nabla_j\nabla_i A^{\ell_1\cdots \ell_q}{}_{k_1\cdots k_p}
\]
\[
= \sum_{r=1}^q R^{\ell_r}{}_{ijm}\,A^{\ell_1\cdots \ell_{r-1}m\ell_{r+1}\cdots \ell_q}{}_{k_1\cdots k_p}
- \sum_{s=1}^p R^{m}{}_{ijk_s}\,A^{\ell_1\cdots \ell_q}{}_{k_1\cdots k_{s-1}mk_{s+1}\cdots k_p}.
\]
\end{remark}
